\documentclass[11pt]{article}

\usepackage[margin=1in]{geometry}
\usepackage[T1]{fontenc}
\usepackage{lmodern}
\usepackage{microtype}
\usepackage{amsmath,amssymb,amsthm,mathtools,bm}
\usepackage{booktabs,longtable,array}
\usepackage{enumitem}
\usepackage[dvipsnames]{xcolor}
\usepackage{hyperref}
\usepackage[nameinlink,capitalise,noabbrev]{cleveref}

\hypersetup{
  colorlinks=true,
  linkcolor=MidnightBlue,
  citecolor=OliveGreen,
  urlcolor=BrickRed,
  pdftitle={The Kannan--Lovasz--Simonovits Conjecture: Problem, Progress, and Research Directions}
}

\setlist[itemize]{topsep=4pt,itemsep=2pt,parsep=1pt}
\setlist[enumerate]{topsep=4pt,itemsep=2pt,parsep=1pt}
\setlength{\parindent}{0pt}
\setlength{\parskip}{5pt}
\allowdisplaybreaks

\newtheorem{theorem}{Theorem}[section]
\newtheorem{proposition}[theorem]{Proposition}
\newtheorem{lemma}[theorem]{Lemma}
\newtheorem{corollary}[theorem]{Corollary}
\theoremstyle{definition}
\newtheorem{definition}[theorem]{Definition}
\newtheorem{target}[theorem]{Open target}
\theoremstyle{remark}
\newtheorem{remark}[theorem]{Remark}

\newcommand{\R}{\mathbb R}
\newcommand{\E}{\mathbb E}
\newcommand{\Var}{\operatorname{Var}}
\newcommand{\Cov}{\operatorname{Cov}}
\newcommand{\Tr}{\operatorname{Tr}}
\newcommand{\op}{\mathrm{op}}
\newcommand{\HS}{\mathrm{HS}}
\newcommand{\Lip}{\mathrm{Lip}}
\newcommand{\Id}{\mathrm{Id}}
\newcommand{\1}{\mathbf 1}
\newcommand{\CP}{C_{\mathrm P}}
\newcommand{\dd}{\,d}
\newcommand{\ip}[2]{\left\langle #1,#2\right\rangle}
\newcommand{\norm}[1]{\left\lVert #1\right\rVert}

\newcommand{\statusbox}[1]{%
  \begin{center}
  \fcolorbox{BrickRed}{BrickRed!5}{%
    \begin{minipage}{0.91\textwidth}
      \textbf{Status.} #1
    \end{minipage}}
  \end{center}}

\title{\textbf{The Kannan--Lov\'asz--Simonovits Conjecture}\\
Problem, Audited Progress, and Possible Future Directions}
\author{Research note from a multi-route proof investigation}
\date{July 2026}

\begin{document}
\maketitle

\begin{abstract}
The Kannan--Lov\'asz--Simonovits (KLS) conjecture asks for a
dimension-free Cheeger, or equivalently Poincar\'e, inequality for every
isotropic log-concave probability measure.  This note gives a standalone
statement of the problem, records the exact reductions needed to pass
between its standard formulations, and summarizes the strongest results
obtained in an extended proof investigation.  The main advances are an
exact second-order ledger for a low eigenfunction, a rank-one-defect
Lichnerowicz theorem, strict spectral amplification under normalized
self-convolution, including its unequal-input extension, and a
posterior reverse-smoothing theorem which makes fixed-Gaussian analytic
outputs a quantitatively equivalent target.  A weighted-Fisher ledger closes for pure translations
and for smooth post-noise Gaussian conditional fibers.  Exact
stochastic-localization and Gaussian-posterior identities isolate the
remaining global gates.  We also record explicit counterexamples to
several tempting forward-only and slice-local claims.  No complete
proof of KLS is claimed.  Checkpoint 17 adds audited posterior
Hilbert--Schmidt and affine-stability ledgers, the asymptotic fixed-Gaussian
saturation comparison, simple-eigenspace and longitudinal Hardy frontiers,
and the product ANOVA/matrix-Stein obstruction.  The final section proposes
a prioritized research program whose statements are precise enough to be
proved or falsified independently.
\end{abstract}

\statusbox{Checkpoint 17 does not contain a complete proof of KLS.
Every assertion labeled ``open target,'' ``candidate,'' or ``future
direction'' remains unproved.  The posterior reverse theorem is a valid
general reduction, while the score and weighted-Fisher estimates close only
near-affine branches.  The affine-interaction inequality needed for the
remaining nonlinear branch is not proved.  All other displayed identities
and model calculations below were checked with their dependence on
dimension and auxiliary parameters exposed.  The new posterior
Hilbert--Schmidt, affine-stability, fixed-Gaussian, simple-eigengap, and
product-ANOVA statements are audited lemmas or conditional frontiers.  The
twice-Bochner ledger still retains an \(O_\kappa(\lambda)\) transverse
saturation remainder and a longitudinal quotient gate; neither is assumed
away.  These results are not a KLS proof.}

\tableofcontents

\section{The problem}

\subsection{Log-concavity, boundary measure, and covariance}

A probability measure \(\mu\) on \(\R^n\) is \emph{log-concave} if for
compact sets \(A,B\subset\R^n\) and \(0<\theta<1\),
\[
 \mu(\theta A+(1-\theta)B)
 \ge \mu(A)^\theta\mu(B)^{1-\theta}.
\]
By Borell's characterization, \(\mu\) is supported on an affine subspace
\(E\); on its minimal affine hull it is either a point mass or has a
density \(e^{-V}\), where \(V\) is convex \cite{Borell1975}.

For a Borel set \(A\), let
\[
 A_\varepsilon=\{x:\operatorname{dist}(x,A)\le\varepsilon\},
 \qquad
 \mu^+(A)=\liminf_{\varepsilon\downarrow0}
 \frac{\mu(A_\varepsilon)-\mu(A)}{\varepsilon}.
\]
The Cheeger constant is
\[
 \psi_\mu=
 \inf_{0<\mu(A)<1}
 \frac{\mu^+(A)}{\min\{\mu(A),1-\mu(A)\}}.
\]
Write \(A_\mu=\Cov(\mu)\).  The measure is \emph{isotropic} when its
barycenter is zero and \(A_\mu=I\).

\subsection{The KLS conjecture}

\begin{theorem}[KLS conjecture; desired statement]
There is a universal constant \(c>0\) such that every non-point
log-concave probability measure \(\mu\) on every Euclidean space satisfies
\[
 \boxed{\qquad
 \psi_\mu\ge \frac{c}{\sqrt{\norm{A_\mu}_{\op}}}.
 \qquad}                                             \tag{KLS}
\]
Equivalently, \(\inf_n\inf_\mu\psi_\mu>0\), where the second infimum is
over isotropic log-concave probabilities on \(\R^n\).
\end{theorem}

The conjecture originated in the study of random walks and isoperimetry
for convex bodies \cite{KLS1995}.  The word \emph{universal} is essential:
no dependence on \(n\), the support, a smoothing parameter, or a
localization time is permitted.

\section{Exact reductions and equivalent targets}

\subsection{Affine hull and degenerate covariance}

Let \(E=m+L\) be the minimal affine hull of the support.  Exterior
Minkowski content computed in \(\R^n\) agrees with that computed using the
induced Euclidean metric on \(E\).  Indeed, for \(B=A\cap E\),
\(B_\varepsilon^E\subset A_\varepsilon\cap E\), while a set already
contained in \(E\) has the same distance function on \(E\) in either
ambient space.  Taking the two infima gives
\[
                         \psi_\mu^{\R^n}=\psi_\mu^E.  \tag{2.1}
\]

On the direction space \(L\), the covariance is positive definite.  If
\(v\in L\) had zero variance, the support would lie in the proper affine
hyperplane \(m+(L\cap v^\perp)\), contradicting minimality.  On
\(L^\perp\), covariance is zero.  Thus the ambient operator norm is the
largest eigenvalue of \(A_\mu|_L\).

\subsection{Affine maps}

If \(S:E\to F\) is invertible and \(\nu=S_\#\mu\), then
\[
 \frac{\psi_\mu}{\norm S_{\op}}
 \le \psi_\nu
 \le \norm{S^{-1}}_{\op}\psi_\mu.                  \tag{2.2}
\]
For example,
\(A_{\varepsilon/\norm S_{\op}}\subset
S^{-1}((SA)_\varepsilon)\), which proves the left inequality directly
from the definition; apply it to \(S^{-1}\) for the other direction.

Translate by \(-m\) and take \(S=A_\mu^{-1/2}\) on \(L\).  The image is
isotropic, and \(\norm{S^{-1}}_{\op}=\sqrt{\norm{A_\mu}_{\op}}\).
Consequently, a universal isotropic lower bound implies (KLS), including
all lower-dimensional supports.  The converse is immediate.  A point mass
is precisely the excluded zero-dimensional case.

\subsection{Poincar\'e and first-moment formulations}

Let \(\CP(\mu)\) be the optimal constant in
\[
 \Var_\mu f\le \CP(\mu)\int|\nabla f|^2\dd\mu.       \tag{2.3}
\]
The chain rule gives
\[
 \CP(S_\#\mu)\le\norm S_{\op}^2\CP(\mu),\qquad
 \CP(\mu)\le\norm{S^{-1}}_{\op}^2\CP(S_\#\mu).    \tag{2.4}
\]
For log-concave measures, Cheeger's inequality and the reverse
Buser--Ledoux implication under \(CD(0,\infty)\), together with standard
convex approximation, give universal constants \(c,C>0\) such that
\[
 \CP(\mu)\le \frac{4}{\psi_\mu^2},
 \qquad
 \psi_\mu\ge\frac{c}{\sqrt{\CP(\mu)}}.              \tag{2.5}
\]
Hence KLS is equivalent, up to constants, to
\[
 \boxed{\quad
 \Var_\mu f\le C\norm{A_\mu}_{\op}
 \int|\nabla f|^2\dd\mu
 \quad}                                               \tag{P}
\]
for every locally Lipschitz \(f\in L^2(\mu)\).  Truncation, cutoff, and
mollification extend a proof for compactly supported smooth tests to this
full class without changing the constant.

E.~Milman's equivalence theorem under the same convexity hypothesis
identifies \(\psi_\mu^{-1}\), up to constants, with
\[
 D_1(\mu)=
 \sup_{\Lip(f)\le1}\int|f-\operatorname{med}_\mu f|\dd\mu
 \cite{Milman2009}.
\]
Mean and median centerings differ by at most a factor two.  Thus the
isotropic form of KLS is also equivalent to
\[
 \sup_{\Lip(f)\le1}\E|f(X)-\E f(X)|\le C.            \tag{M}
\]
Linear functions already satisfy this by isotropy.  Radial Lipschitz
functions are controlled by the dimension-free thin-shell theorem.  A bad
witness must therefore be genuinely nonlinear and non-radial.

\begin{remark}
KLS is not a dimension-free log-Sobolev statement.  Isotropic
one-dimensional exponential measures already obstruct subgaussian tails.
No argument below assumes a uniform log-Sobolev or \(T_2\) inequality.
\end{remark}

\section{An exact low-eigenfunction ledger}

This section records the most useful general identities found in the
investigation.  Work first with a smooth convex potential on a smooth
convex support and Neumann boundary condition; the closed-form identities
extend by convex approximation.

Let \(A=-L_\mu\) be the nonnegative weighted generator and suppose
\[
 Af=\lambda f,\qquad
 \E f=0,\qquad \E f^2=1,\qquad \E|\nabla f|^2=\lambda.
\]
Set
\[
 a=\E[Xf],\qquad r=f-a\cdot x,\qquad \delta=\norm r_2.
\]

\begin{proposition}[Affine residual identities]
For an isotropic \(\mu\),
\begin{align}
 &\E r=0,\qquad \E[Xr]=0,
 \qquad |a|^2=1-\delta^2,                              \tag{3.1}\\
 &\E\nabla f=\lambda a,
 \qquad \E\nabla r=-(1-\lambda)a,                    \tag{3.2}\\
 &\E|\nabla r|^2
   =1-\lambda+(2\lambda-1)\delta^2,                  \tag{3.3}\\
 &\Var(\nabla f)
   =\lambda\bigl[1-\lambda(1-\delta^2)\bigr].       \tag{3.4}
\end{align}
\end{proposition}

\begin{proof}
The first line is orthogonal projection onto the linear functions.
Testing the weak eigenvalue equation against \(b\cdot x\) gives
\(b\cdot\E\nabla f=\lambda b\cdot a\), proving (3.2).
The remaining formulas follow by expanding
\(\nabla r=\nabla f-a\).
\end{proof}

Let
\[
 B_f=\E\norm{D^2f}_{\HS}^2
\]
and let \(\mathcal R_V(f)\ge0\) denote the sum of the potential-curvature
and Reilly boundary-curvature terms.  Bochner--Reilly gives
\[
 B_f+\mathcal R_V(f)=\lambda^2.                       \tag{3.5}
\]
Applying the same bottom Poincar\'e inequality componentwise to
\(\partial_i f\) and using (3.4) produces the exact nonnegative defect
identity
\[
 \boxed{\quad
 [B_f-\lambda\Var(\nabla f)]+\mathcal R_V(f)
 =\lambda^3(1-\delta^2).
 \quad}                                               \tag{3.6}
\]
In particular,
\[
                         \norm{D^2r}_2\le\lambda.     \tag{3.7}
\]

The identity says that a hypothetical small-gap, near-linear eigenfunction
has a residual \(r\) that is small both in value and Hessian, yet has mean
gradient almost \(-a\).  This isolates the following gate.

\begin{target}[Residual mean-gradient inequality]
Prove, for every centered isotropic log-concave \(\mu\),
\[
 \boxed{\qquad
 |\E\nabla g|
 \le C\bigl(\norm g_2+\norm{D^2g}_2\bigr),
 \quad \E g=0,\quad \E[Xg]=0.
 \qquad}                                              \tag{MG}
\]
\end{target}

Applied to \(r\), (MG) gives
\[
 (1-\lambda)\sqrt{1-\delta^2}\le C(\delta+\lambda), \tag{3.8}
\]
and therefore closes the near-linear branch whenever \(\delta\) is a
sufficiently small universal constant.

The stronger Hessian-only form holds in one dimension and tensorizes to
products.  If \(\nu\) is centered, variance one, and log-concave on
\(\R\), its one-dimensional Poincar\'e constant is universally bounded.
Applying Poincar\'e first to \(h'\) and then to \(h-(\int h'\dd\nu)x\)
gives, for \(h\perp\{1,x\}\),
\[
 \left|\int h'\dd\nu\right|
 \le12\left(\int(h'')^2\dd\nu\right)^{1/2}.          \tag{3.9}
\]
Conditional projections onto product coordinates are orthogonal, so
\[
 |\E\nabla g|^2
 \le144\E\norm{D^2g}_{\HS}^2                         \tag{3.10}
\]
for products of isotropic one-dimensional log-concave laws.

\subsection{A false interpolation and its diagnostic value}

The seemingly natural inequality
\[
 \Var(\nabla f)\le C\delta\norm{D^2f}_2              \tag{3.11}
\]
is false even for genuine first eigenfunctions.  The first degree-one
Neumann mode on the isotropic Euclidean ball has
\[
 \delta\asymp n^{-1},\qquad
 \Var(\nabla f)\asymp n^{-1},\qquad
 \norm{D^2f}_2\asymp n^{-1/2},
\]
so the ratio in (3.11) grows like \(\sqrt n\).  Any viable near-linearity
argument must target the \emph{mean} gradient or another finite-dimensional
component, not the whole gradient variance.

\section{The directional third-moment gate}

For a symmetric matrix \(B\), define
\[
 q_B=X^TBX-\Tr B,
 \qquad
 T B=\E[Xq_B].
\]
The quadratic restriction of the Hessian-only form of (MG) is exactly
\[
 \boxed{\qquad
 \sup_{|u|=1}
 \norm{\E[(u\cdot X)(XX^T-I)]}_{\HS}\le C.
 \qquad}                                              \tag{TM}
\]
Indeed, \(T^*u=\E[(u\cdot X)(XX^T-I)]\), while
\(D^2q_B=2B\).  More explicitly, apply the inequality to
\(q_B-(TB)\cdot X\), which is orthogonal to both constants and linear
functions and has the same Hessian as \(q_B\).

Thin-shell estimates for every projection give only
\[
 |\Tr(PM_u)|\le C\sqrt{\operatorname{rank}P},
 \qquad M_u=\E[(u\cdot X)(XX^T-I)].                  \tag{4.1}
\]
Ordering the eigenvalues yields
\(|\lambda_j(M_u)|\le Cj^{-1/2}\), hence
\[
                         \norm{M_u}_{\HS}\le C\sqrt{\log(en)}. \tag{4.2}
\]
The harmonic diagonal matrix shows that (4.1) alone cannot remove the
logarithm.  A successful proof of (TM) must use compatibility among matrix
directions, not only rank-by-rank trace information.

\subsection{Dependent model classes where (TM) is dimension-free}

The investigation established (TM) with universal constants for several
nonproduct classes.

\paragraph{Paired wedges.}
For
\[
 K_a=\{(t,x,y):|t|\le1,
 |x_i|\le1+a_it,
 |y_i|\le1-a_it\},
\]
let \(S=\sum_i a_i^2\) and \(v=\Var T\).  The axial marginal potential
satisfies \(W''\ge2S\), hence \(vS\le1/2\).  After whitening, the only
nonzero third coefficients are
\[
 \lambda_i=\frac{2a_i\sqrt v}{1+a_i^2v},
 \qquad \sum_i\lambda_i^2\le2,
\]
and therefore
\[
 \sup_{|u|=1}\norm{M_u}_{\HS}\le2.                  \tag{4.3}
\]
The whitened law is an explicit \(4\)-Lipschitz image of a product law,
which also gives \(\CP\le192\).  A one-sided wedge admits an analogous
map and the explicit bound \(\CP<972\).

\paragraph{Dirichlet laws.}
For a log-concave Dirichlet distribution with parameters
\(\alpha_i\ge1\), write \(A=\sum_i\alpha_i\) and
\(q_i=\alpha_i/A\).  A unit tangent direction has coefficients \(a_i\)
with \(\sum q_ia_i=0\) and \(\sum q_ia_i^2=1\).  Direct moment algebra gives
\[
 \norm{M_{u(a)}}_{\HS}^2
 =\frac{4(A+1)}{(A+2)^2}
   \left(\sum_i a_i^2-2\right)<4.                    \tag{4.4}
\]

\paragraph{Cones and conditional Gaussian families.}
For a cone over a centrally symmetric base,
\[
 \sup_{|u|=1}\norm{M_u}_{\HS}^2
 =\frac{4(n-1)(n+2)}{(n+3)^2}<4.                     \tag{4.5}
\]
The same gate holds for affine box slices and affine Gaussian covariance
pencils.

More generally, consider
\[
 p(t,y)\propto
 \exp\!\left[-W(t)-\frac12y^TQ(t)y\right],
 \qquad R(t)=Q(t)^{-1}.
\]
Joint log-concavity is equivalent to
\[
                         W''\ge0,\qquad R''\preceq0.  \tag{4.6}
\]
This permits genuinely rotating, noncommuting covariance eigenspaces.
After whitening, marginal curvature controls the normalized covariance
velocity by
\[
 \Phi''(t)\ge\frac12
 \norm{R(t)^{-1/2}R'(t)R(t)^{-1/2}}_{\HS}^2.          \tag{4.7}
\]
A one-dimensional Stein-kernel argument then proves (TM) with a universal
constant for the entire conditional-Gaussian class.

\subsection{Why general slicing remains open}

For an arbitrary smooth conditional law
\(p_t(y)\propto e^{-V(t,y)}\), set
\(H=D_y^2V\), \(g=V_t-\E_tV_t\), and let \(Lu=g\) for the conditional
generator.  The Brascamp--Lieb deficit controls
\(\E_t\norm{D^2u}_{\HS}^2\).  Covariance velocity would follow from
\[
 \boxed{\quad
 \norm{\E[Y\otimes\nabla u+\nabla u\otimes Y]}_{\HS}^2
 \le C\E\norm{D^2u}_{\HS}^2.
 \quad}                                               \tag{AH}
\]
This affine-Hessian rigidity is itself a near-linear KLS gate.  Moving
boundaries cannot be omitted: for an interval and for a centered
exponential law, the covariance of the iid sum-difference slice is convex,
not concave, and all relevant curvature is carried by endpoint flux.

\section{A rank-one-defect Lichnerowicz theorem}

One genuinely general positive theorem arose from the localization route.

\begin{proposition}[Rank-one curvature defect]
Let \(u\in S^{n-1}\), \(P=I-u\otimes u\), and let
\(d\mu=e^{-V}dx\) be log-concave.  If, distributionally,
\[
                         D^2V\succeq\kappa P,
\]
then
\[
 \boxed{\quad
 \CP(\mu)\le201\left(\kappa^{-1}
 +\Var_\mu\ip{X}{u}\right).
 \quad}                                               \tag{5.1}
\]
The same statement holds intrinsically on a lower-dimensional affine
support.
\end{proposition}

Here is the proof mechanism.  Disintegrate at
\(x=tu+y\), with conditional law \(\nu_t\) on \(u^\perp\) and marginal
\(e^{-U(t)}dt\).  There is a canonical continuity-equation velocity
\(w_t\) satisfying
\[
 \partial_t\nu_t+\operatorname{div}(\nu_tw_t)=0,
 \qquad
 \kappa\int|w_t|^2\dd\nu_t\le U''(t).                \tag{5.2}
\]
The second inequality is a reinforced infinitesimal Pr\'ekopa theorem;
the cross-Hessian block is essential.  If
\(g(t)=\int f(t,y)\dd\nu_t(y)\), then
\[
 g'(t)=\int(\partial_tf+\nabla_yf\cdot w_t)\dd\nu_t,
\]
so
\[
 \frac{g'(t)^2}{1+\int|w_t|^2\dd\nu_t}
 \le\int|\nabla f|^2\dd\nu_t.                        \tag{5.3}
\]
A curvature-weighted one-dimensional Hardy inequality, plus conditional
Lichnerowicz in \(u^\perp\), proves (5.1).  Gaussian convolution preserves
the transverse curvature quantitatively:
\[
 \kappa\longmapsto\frac{\kappa}{1+\kappa\varepsilon},
\]
which permits passage to nonsmooth potentials without loss in the final
constant.

This theorem is useful after localization has created curvature in all but
one direction.  Its remaining obstruction is the variance of the
adaptively selected survivor direction.

\section{Stochastic localization and heat-flow identities}

\subsection{Function-specific localization}

For a normalized first eigenfunction, a localization driver can preserve
its posterior mean while adding curvature away from one selected direction.
Combining the preserved variance with (5.1) gives the exact ledger
\[
 \boxed{\quad
 1\le\frac{192}{T}\lambda
 +96\E[R_T\mathcal E_T].
 \quad}                                               \tag{6.1}
\]
Here \(R_T\) is the terminal variance in the surviving direction and
\(\mathcal E_T=\E_T|\nabla f|^2\) is posterior energy.  Thus a small gap
forces \(R_T\) to be of order \(1/\lambda\) under the
energy-biased path law.  Ordinary expectation bounds on \(R_T\) are
insufficient because of this correlation.

A softened driver that creates full curvature has the complementary
tradeoff
\[
                         1-\E m_T^2\le\frac{\lambda}{\eta T}, \tag{6.2}
\]
where \(m_T\) is the localized mean and \(\eta T I\) is the curvature
created.  Equations (6.1)--(6.2) make precise why preserving the signal and
regularizing its selected direction compete.

\subsection{Parallel coupling and a low-mode martingale}

In standard stochastic localization, write \(A_t=\Cov(\mu_t)\) and let
\[
 M_0=I,\qquad \dot M_t=A_tM_t.
\]
For a fixed \(f\), set
\[
 c_t=\Cov_{\mu_t}(X,f),\qquad
 D_t=\E_t[(f-\E_tf)(X-a_t)(X-a_t)^T].
\]
Direct It\^o differentiation gives
\[
 dc_t=D_t\dd B_t-A_tc_t\dd t,
 \qquad
 \boxed{\ M_t^Tc_t\text{ is a martingale}.\ }        \tag{6.3}
\]
For the eigenfunction decomposition in \cref{sec:posterior} below,
\(c_0=a\).  At time one, \(\mu_1\) is 1-strongly log-concave, so
\[
                         \E|c_1|^2\le\lambda.         \tag{6.4}
\]
With \(b=a/|a|\), (6.3)--(6.4) imply
\[
 \boxed{\qquad
 \E|M_1b|^2\ge\frac{1-\delta^2}{\lambda}.
 \qquad}                                              \tag{6.5}
\]
A nearly linear small mode must therefore select an expensive direction
of the coupling.  Available theorems control only the trace
\[
                         \E\Tr(M_1^TM_1)\le Cn,        \tag{6.6}
\]
not the selected direction.

The weakest final-time statement that closes this branch is
\[
 \boxed{\quad
 \E\left[
 \frac{\ip{M_1b}{c_1}^2}{|c_1|^2}\1_{\{c_1\ne0\}}
 \right]\le C.
 \quad}                                               \tag{FA}
\]
This is strictly weaker than \(\E|M_1b|^2\le C\): it controls only the
random normal component used by the martingale pairing.  It passes the
Gaussian, cube, product-exponential, ball, and simplex tests.  A one-spike
matrix path shows that (6.6) and ordered-rank stopping bounds alone do not
imply (FA).

Coarea does not automatically help.  The quotient in (FA) is a normal
stretch of the initial-tilt flow at one fixed starting point.  Change of
variables averages over random preimages instead, and endpoint integration
by parts controls the conditional mean of \(M_1b\), not its conditional
second moment.

\section{Gaussian posterior and a directional-MMSE target}
\label{sec:posterior}

Let \(X\sim\mu\), \(G\sim N(0,I)\) be independent, and \(Y=X+G\).
Write
\[
 m(y)=\E[X\mid Y=y],\qquad
 A(y)=\Cov(X\mid Y=y).
\]
The posterior is 1-strongly log-concave, hence \(0\preceq A(y)\preceq I\).
For
\(F(y)=\E[f(X)\mid Y=y]\), exponential-family differentiation gives
\[
                         \nabla F(y)=\Cov(X,f\mid Y=y)=c(y). \tag{7.1}
\]
If \(f=a\cdot x+r\), then exactly
\[
                         c(y)=A(y)a+e(y),
 \qquad e(y)=\Cov(X,r\mid Y=y).                       \tag{7.2}
\]
Posterior Poincar\'e and total variance imply
\[
 \E|c(Y)|^2\le\lambda,
 \qquad
 \E|e(Y)|^2\le\delta^2.                              \tag{7.3}
\]

For a unit vector \(b\), define
\[
 \operatorname{mmse}_\mu(b)
 =\E\Var(b\cdot X\mid X+G)
 =\E[b^TA(Y)b].                                       \tag{7.4}
\]

\begin{target}[Directional MMSE floor]
Prove that every centered isotropic log-concave law satisfies
\[
 \boxed{\qquad
 \inf_{|b|=1}\operatorname{mmse}_\mu(b)\ge c_0>0.
 \qquad}                                              \tag{DMMSE}
\]
\end{target}

This would close the very-near-linear branch immediately.  Indeed, with
\(a=|a|b\), (7.2)--(7.3) and the Hilbert-space triangle inequality give
\[
 \sqrt\lambda
 \ge c_0\sqrt{1-\delta^2}-\delta.                    \tag{7.5}
\]

The target is nonlinear despite involving a linear latent variable.  The
orthogonal noisy coordinates may carry nonlinear information about
\(b\cdot X\).  Tweedie's identity
\[
 m(y)=y+\nabla\log q(y),\qquad \nabla m(y)=A(y),       \tag{7.6}
\]
where \(q\) is the density of \(Y\), shows the KLS content.  For
\(h_b(y)=b\cdot m(y)\),
\[
 \Var_q(h_b)=1-\operatorname{mmse}_\mu(b),
 \qquad
 \int|\nabla h_b|^2\dd q=\E|A(Y)b|^2.                \tag{7.7}
\]
Thus failure of (DMMSE) produces a low-energy monotone posterior-mean
function for the log-concave Gaussian convolution \(q\).

There is a possible dimension-descent estimate worth auditing.  Reveal
\(Y_\perp=P_{b^\perp}Y\) first and put
\[
 v(z)=\Var(b\cdot X\mid Y_\perp=z),\qquad
 h_0(z)=\E[b\cdot X\mid Y_\perp=z].
\]
The rank-one curvature defect gives transverse posterior covariance at
most \(I\), hence
\[
                         |\nabla h_0(z)|^2\le v(z).    \tag{7.8}
\]
Conditional one-dimensional log-concavity gives a scalar Gaussian-channel
MMSE comparable from below to \(\min(v,1)\).  One-dimensional log-concave
moments suggest
\(\norm v_p\lesssim p^2\).  Combining these facts with Poincar\'e for the
\((n-1)\)-dimensional marginal suggests the candidate recursion
\[
 \operatorname{mmse}_\mu(b)
 \gtrsim
 \frac{1}{\CP(q_\perp)\log^2(e+\CP(q_\perp))}.        \tag{7.9}
\]
Equation (7.9) was not fully audited before the investigation stopped.  It
would still contain logarithmic loss, but it offers a concrete starting
point for a new dimension-descent mechanism.

\section{Further exact progress}

\subsection{Convex half-set distance}

If \(A\) is convex and \(\mu(A)=1/2\), one-dimensional localization of
parallel expansions gives the dimension-free lower bound
\[
 \boxed{\qquad
 \E\operatorname{dist}(X,A)
 \ge \frac{2\log2-1}{4\mu^+(A)}.
 \qquad}                                              \tag{8.1}
\]
The scalar constant is explicit.  The reverse inequality for arbitrary
half-sets is already KLS-strength, and naive convexification fails even for
the isotropic Laplace law.

\subsection{Tensor slow-growth plateaus}

Let
\[
 \mathcal C_n=\sup\{\CP(\mu):\mu\text{ isotropic log-concave on }\R^n\}.
\]
If \(\mathcal C_n\) were unbounded while obeying a logarithmic upper
envelope, monotonicity implies dimensions \(n_j\) and integers
\(r_j\to\infty\) such that
\[
 \frac{\mathcal C_{r_jn_j}}{\mathcal C_{n_j}}\to1.   \tag{8.2}
\]
Thus large tensor powers are relative near-extremizers with growing first
eigenspace multiplicity.

This does not by itself create a contradiction.  For a mixed tensor mode
indexed by a symmetric zero-diagonal matrix \(B\), the spectral splitting
is at most the convex-buffer cost times
\[
                         \frac{\norm B_{\op}}{\norm B_{\HS}}\le1. \tag{8.3}
\]
Moreover, the \(r\) tensor low modes occupy a rank-\(r\) exceptional
subspace in ambient dimension \(rn\); ordered-rank localization still sees
\(\log((rn)/r)=\log n\).  Multiplicity alone therefore does not remove the
small-rank obstruction.

\section{Self-convolution and the weighted-slice program}

\subsection{Strict forward amplification under self-convolution}

Let \(X,Y\) be independent with the same centered isotropic log-concave
law \(\mu\), and put
\[
 S=\frac{X+Y}{\sqrt2},\qquad \nu=\mathcal L(S),\qquad
 a=\lambda_1(\mu)=\CP(\mu)^{-1},\qquad
 b=\lambda_1(\nu)=\CP(\nu)^{-1}.
\]

\begin{proposition}[Forward self-convolution amplification]
Without assuming that either spectral edge is attained,
\[
 \boxed{\qquad
 b(1-b)\le4(b-a),\qquad
 b\ge\frac{\sqrt{9+16a}-3}{2}.
 \qquad}                                              \tag{SC}
\]
\end{proposition}

For a centered unit test \(f(S)\), set
\[
 h(X)=\E[f(S)\mid X],\quad
 R_f=f(S)-h(X)-h(Y),\quad q=\E|\nabla f(S)|^2.
\]
The product Poincar\'e inequality on the two-coordinate-centered residual
and the vector Hoeffding decomposition give
\[
 E_R:=\E|\nabla_{X,Y}R_f|^2\le2(q-a),                \tag{10.1}
\]
\[
 \Var(Z)=2\Var(U)+\E|W|^2,\qquad
 E_R=\Var(U)+\E|W|^2,                                \tag{10.2}
\]
where \(Z=\nabla f(S)\), \(U=\E[Z\mid X]\), and
\(W=Z-\E[Z\mid X]-\E[Z\mid Y]+\E Z\).  Thus
\(\Var(\nabla f)\le2E_R\le4(q-a)\).  Taking \(f\) in a
spectral window \([b,b+\varepsilon]\), testing the weak generator equation
against coordinates, and sending \(\varepsilon\downarrow0\) gives
\(\Var(\nabla f)\ge b(1-b)-o(1)\).  This proves (SC), including at a
continuous bottom spectral edge.

If \(a_k\) is the gap after \(k\) normalized dyadic convolutions, then
\(a_k\uparrow1\), and
\[
 a_{k+1}\le\frac12\quad\Longrightarrow\quad
 a_{k+1}\ge\frac87a_k.                               \tag{10.3}
\]
This is only forward amplification.  It may take
\(O(\log(1/a_0))\) convolutions to reach a fixed gap, so (SC) gives no
dimension-free lower bound on the initial \(a_0\).

\subsection{Unequal inputs and fixed-Gaussian output structure}

Equality of the two input laws is not needed.  Let \(\mu_0,\mu_1\) be
centered isotropic log-concave laws on the same Euclidean space, let
\(X\sim\mu_0\), \(Y\sim\mu_1\) be independent, and set
\[
 S=\frac{X+Y}{\sqrt2},\qquad
 a=\min\{\lambda_1(\mu_0),\lambda_1(\mu_1)\},\qquad
 b=\lambda_1(\mathcal L(S)).
\]
The unequal vector predictors need not have the same variance.  If
\[
 u=\E|\E[\nabla f(S)\mid X]-\E\nabla f(S)|^2,\quad
 v=\E|\E[\nabla f(S)\mid Y]-\E\nabla f(S)|^2
\]
and \(w\) is the squared norm of the two-coordinate residual, then
\[
 \Var(\nabla f)=u+v+w,\qquad
 E_R=\frac{u+v}{2}+w.
\]
Consequently the proof of (SC) is unchanged:
\[
 \boxed{\qquad
 b(1-b)\le4(b-a),\qquad
 b\ge\frac{\sqrt{9+16a}-3}{2}.
 \qquad}                                              \tag{USC}
\]
The identities are first read for smooth tests and then in the closed
weighted Sobolev form domain.  For log-concave laws, smooth tests form a
core after working intrinsically on the minimal affine hull.  This note
does not extend (USC) to arbitrary singular non-log-concave inputs.

Now take \(\mu_1=\gamma=N(0,I)\) and define
\[
 \mathcal G\mu
 =\mathcal L\left(\frac{X+G}{\sqrt2}\right).
\]
Since every isotropic law has
\(\lambda_1(\mu)\le1=\lambda_1(\gamma)\), the input parameter in
(USC) is \(a=\lambda_1(\mu)\).  The only rearrangement supplied by
(USC) is
\[
 \lambda_1(\mu)
 \le\frac{\lambda_1(\mathcal G\mu)
 (3+\lambda_1(\mathcal G\mu))}{4}.                   \tag{10.3a}
\]
This forward constraint gives no lower bound on the input from a lower
bound on the output.  Hence (USC) alone gives no fixed-Gaussian reduction
of KLS.  There is, however, a separate reverse comparison.

\begin{proposition}[Posterior reverse smoothing]
Let
\(a=\lambda_1(\mu)\) and
\(b=\lambda_1(\mathcal G\mu)\).  Then
\[
 \boxed{\qquad
 b\le\frac{2a}{1-a}\quad(a<1),
 \qquad a\ge\frac{b}{2+b}.
 \qquad}                                             \tag{10.3c}
\]
The second inequality also holds when \(a=1\).
\end{proposition}

Indeed, take a centered normalized input form-domain test \(f\) with
energy \(q<1\), put \(Y=X+G\), and set
\(F(y)=\E[f(X)\mid Y=y]\).  The posterior is intrinsically
one-strongly log-concave, so
\[
 \E\Var(f(X)\mid Y)\le q,
 \qquad
 \Var(F(Y))\ge1-q.                                  \tag{10.3d}
\]
Posterior differentiation gives
\(\nabla F=\Cov(X,f\mid Y)\).  Covariance
Cauchy--Schwarz together with
\(\Cov(X\mid Y)\preceq I\) gives the dimension-free pointwise estimate
\[
 |\nabla F(Y)|^2\le\Var(f(X)\mid Y),
\]
and hence the energy of \(s\mapsto F(\sqrt2s)\) under
\(\mathcal G\mu\) is at most \(2q\).  The variational definition of
\(b\), followed by a form-domain minimizing sequence \(q\downarrow a\),
proves (10.3c).  The argument is intrinsic on a lower-dimensional affine
support and uses posterior strong log-concavity for hard supports, so no
spectral attainment or ambient full-dimensionality is assumed.

The output has a positive analytic log-concave density.  Before the
\(1/\sqrt2\) rescaling, if \(U=-\log(\mu*\gamma)\), the Gaussian posterior
identity and posterior Brascamp--Lieb inequality give
\[
 D^2U(y)=I-\Cov(X\mid X+G=y),\qquad
 0\preceq\Cov(X\mid X+G=y)\preceq I.
\]
For the potential \(\widetilde U(s)=U(\sqrt2s)+\mathrm{const}\) of
\(\mathcal G\mu\),
\[
                         0\preceq D^2\widetilde U\preceq2I. \tag{10.3b}
\]
The forward inequality (USC) does not reverse, but the independent
posterior estimate (10.3c) does.  Therefore the positive analytic isotropic
class
\[
                         0\preceq D^2\widetilde U\preceq2I
\]
is a quantitatively equivalent KLS target: a universal output gap
\(b\ge c_0\) gives \(a\ge c_0/(2+c_0)\), while the converse direction is
restriction to a subclass.  This upper-curvature condition is not a
Bakry--\'Emery lower-curvature hypothesis.  Lower-dimensional inputs and
their Gaussian noises are handled intrinsically on their affine hulls.

\subsection{Score rigidity, affine interaction, and one-form circularity}

On the regularized class \(0\preceq D^2V\preceq\beta I\), integration by
parts and covariance positivity give the score-rigidity estimate
\[
 |\E\nabla g|\le\sqrt{\beta-1}\,\norm g_2           \tag{10.3e}
\]
for every affine-orthogonal \(g\).  At \(\beta=2\), this gives a numerical
gap for every bottom spectral-window vector whose \(L^2\) affine residual
is at most \(1/2\).  The unresolved nonlinear branch can be isolated by
the affine-interaction candidate
\[
 \operatorname{dist}_{L^2(\mathcal G\mu)}(f,\mathrm{Aff})^2
 \le C_A\E\left|f\left(\frac{X+G}{\sqrt2}\right)
 -\E[f\mid X]-\E[f\mid G]\right|^2.                 \tag{AI}
\]
Its kernel is exactly affine in full dimension and its sharp Gaussian
constant is two, but no universal proof is supplied here; already its
quadratic sector contains a dimension-free generalized quadratic-variance
theorem.

The tempting generic replacement
\[
 \int|w|^2d\nu\le C\left\{\int\|Dw\|_{\HS}^2d\nu
              +\int w^TD^2Vw\,d\nu\right\},
 \qquad w=\nabla f,                                  \tag{10.3f}
\]
does not provide an extra lemma.  Bochner and spectral calculus show that
its optimal constant is exactly \(C_P(\nu)\); the centered-gradient Korn
constant likewise lies between \(C_P(\nu)-1\) and \(C_P(\nu)\).
Thus (10.3f) is circular.  The actual remaining mechanism is coherence of
an almost-centered, slowly varying field which adaptively saturates the
posterior covariance, or a proof of (AI) on the relevant bottom spectral
windows.

\subsection{Checkpoint 17: posterior saturation and affine stability}

Let \(\pi\) be centered \(1\)-strongly log-concave on its intrinsic
Euclidean support, let \(Z\sim\pi\), let \(A=\Cov Z\), and put
\(\delta_\pi(u)=|u|^2-u^TAu\).  The audited mixed-third-moment theorem is
\[
 M_\pi(u):=\E[(u\cdot Z)ZZ^T],\qquad
 \norm{M_\pi(u)}_{\HS}\le4\sqrt{\delta_\pi(u)}.       \tag{17.1}
\]
For \(Y=X+G\) with intrinsic \(G\sim N(0,I)\), write
\(m_y=\E[X\mid Y=y]\), \(Z_y=X-m_y\),
\(A_y=\Cov(X\mid Y=y)\), and \(H(y)=I-A_y\).  Differentiation gives
\(D A_y[v]=\E[(v\cdot Z_y)Z_yZ_y^T\mid Y=y]\), and hence, for every
(even adaptively selected) \(u=u(y)\),
\[
 \left(\sum_j\|(D A_y[e_j])u\|^2\right)^{1/2}
 \le4\sqrt{u^TH(y)u}.                                \tag{17.2}
\]
This is a basis-summed Hilbert--Schmidt estimate with no trace or rank
factor; it is local and does not synchronize directions globally.

For every locally Sobolev vector field \(W\), the same estimate and
\(0\preceq H\preceq I\) give
\[
 D(HW)[v]=(D H[v])W+H(DW[v]),\qquad
 \norm{D(HW)}_{\HS}\le\norm{DW}_{\HS}+4\sqrt{W^THW},
\]
and therefore, with
\(\mathcal D=\E\norm{DW}_{\HS}^2\) and
\(\mathcal K=\E[W^THW]\),
\[
 \E\norm{D(HW)}_{\HS}^2\le2\mathcal D+32\mathcal K,
 \qquad \E|HW|^2\le\mathcal K.                     \tag{17.3}
\]

The audited affine-stability lemma says that for centered
\(1\)-strongly log-concave \(\pi\),
\[
 D_\pi(f):=\E|\nabla f|^2-\Var_\pi f\ge0
 \quad\Longrightarrow\quad
 \inf_{c,\ell}\E|f(Z)-c-\ell\cdot Z|^2\le22D_\pi(f). \tag{17.4}
\]
For a bottom input spectral-window vector \(g\) in the channel \(Y=X+G\),
put \(q=\E|\nabla g|^2\) and
\(d=\E\Var(g(X)\mid Y)\).  The reverse comparison gives
\(q-d=O(a^2)\) whenever \(a=\lambda_1(\mu)\downarrow0\) (window width
\(o(a^2)\)).  Applying (17.4) in each posterior gives
\[
 \E_Y\inf_{c,\ell}\E_{\pi_Y}
 |g(X)-c-\ell\cdot(X-m_Y)|^2
 \le22(q-d)=O(a^2).                                 \tag{17.5}
\]
The affine coefficients may switch with \(y\); (17.5) is not a global
Gaussian-factor theorem.

For the same intrinsic channel, let
\(\lambda=\lambda_1(\mathcal L(Y))\) and let \(b=2\lambda\) be the gap of
the normalized output.  Spectral-window passage gives
\[
 \frac{\sqrt{1+12a}-1}{6}\le\lambda\le\frac{a}{1-a},
 \qquad a-3a^2\le\lambda\le a+2a^2\quad(a\downarrow0),
\]
and therefore
\[
 \lambda=a+O(a^2),\qquad b=2a+O(a^2).                 \tag{17.6}
\]
No spectral-edge attainment is used in these comparisons.

\subsection{Simple eigengap and the longitudinal frontier}

Let \(H=D^2U\), let \(h(y)\) be its simple lowest eigenvalue, and assume
\(\lambda_2(H(y))-h(y)\ge\kappa>0\).  If \(P(y)\) is the rank-one lowest
spectral projection, the audited Davis--Kahan/Hilbert--Schmidt calculation
gives
\[
 \sum_j\norm{D_jP(y)}_{\HS}^2\le \frac{32}{\kappa^2}h(y). \tag{17.7}
\]
Writing \(Q=I-P\), the field ledger (with \(\mathcal D,\mathcal K\) as
above) is
\[
 \E|QW|^2\le\mathcal K/\kappa,\qquad
 \E\norm{D(PW)}_{\HS}^2
 \le2\mathcal D+\frac{32}{\kappa^2}\mathcal K.       \tag{17.8}
\]

The twice-Bochner ledger is dimension-free but leaves an additional
transverse remainder.  With the Witten one-form operator
\(\mathcal A_1=-L+H\) and \(C=DW=\alpha P+C_0\),
\(\alpha=e^TCe\), and \(t=D^3U[e,e,W]\), it reads
\[
 \begin{aligned}
 \norm{\mathcal A_1W}_2^2\ge{}&
 \E\norm{D(DW)}_{\HS}^2+3\E[h\alpha^2]
 +\frac{3\kappa}{4}\E\norm{C_0}_{\HS}^2+2\E[\alpha t]+\E|HW|^2\\
 &\hspace{18mm}-\frac{64}{3\kappa}\mathcal K .
 \end{aligned}                                         \tag{17.9}
\]
For a normalized bottom field the final term is \(O_\kappa(\lambda)\), not
\(O_\kappa(\lambda^2)\), so it is not absorbed by (17.9).  The remaining
frontier therefore has both this transverse saturation remainder and the
scalar longitudinal quotient gate.

If the low line is fixed, the latter gate closes by the one-dimensional
Hardy factorization.  For a one-dimensional log-concave \(q=e^{-u}\) with
\(\Var_q(s)\le2\), and \(h=u''\), every smooth compactly supported amplitude
\(a\) satisfies
\[
 \E(a'')^2+3\E[h(a')^2]+2\E[aa'h']+\E[h^2a^2]
 \ge \frac1{24}\{\E(a')^2+\E[ha^2]\}.                \tag{17.10}
\]
Thus a nonzero scalar one-form eigenmode has eigenvalue at least \(1/24\),
and the fixed-global-eigenline/product branch is closed.  A generic
amplitude-connectivity replacement is false even for exact analytic
gradients: for \(\nu=N(0,2I_2)\),
\(f_R=(e^{Rx}+e^{Ry})/R\), and \(W_R=\nabla f_R\), define
\(d\sigma_R=|W_R|^2d\nu/(2e^{4R^2})\) and
\(N_R=W_RW_R^T/|W_R|^2\).  Then
\[
 \inf_{\operatorname{rank} P_0=1}\E_{\sigma_R}\norm{N_R-P_0}_{\HS}^2
 =1-e^{-2R^2},\qquad
 \E_{\sigma_R}\norm{DN_R}_{\HS}^2\le R^2e^{-2R^2}.  \tag{17.11}
\]
This example has repeated \(H\) and high field energy, so it does not
disprove the simple-low-mode program; it does disprove generic
amplitude-biased Poincare/connectivity claims.

\subsection{Exact ANOVA jet, product closure, and matrix-Stein obstruction}

For \(\mu=\bigotimes_i\mu_i\), with every factor centered,
variance-one, and log-concave, let \(f\in L^2(\mathcal L(S))\) be
centered and put \(S=(X+G)/\sqrt2\), \(h(X)=\E[f(S)\mid X]\), and
\(R=f(S)-\E[f\mid X]-\E[f\mid G]\), Gaussian Hermite jets give
\[
 \E R^2=\sum_{|\alpha|\ge1}\frac{\Var_\mu(\partial^\alpha h)}{\alpha!},
 \qquad
 \norm{f}_2^2=V_0+\sum_{k\ge1}(V_k+M_k).
\]
The diagonal one-dimensional Stein ledger gives the dimension-free product
affine-interaction theorem
\[
 \operatorname{dist}_{L^2(\mathcal L(S))}(f,\mathrm{Aff})^2\le900\,\E R^2. \tag{17.12}
\]
For a general centered isotropic law, a conditional frame-free criterion is
the following: if a matrix Stein kernel \(\tau\) satisfies \(\E\tau=I\) and
\[
 K^2=\norm{\E[(\tau-I)(\tau-I)^T]}_{\op}<\infty,
\]
then, with Poincare constant \(C_0\), the same jet argument gives
\[
 \operatorname{dist}(f,\mathrm{Aff})^2\le(C_0+3+2K^2)\,\E R^2.             \tag{17.13}
\]
Known general log-concave kernels have only a Hilbert--Schmidt/trace budget
of order \(n\); the operator bound in (17.13) is unresolved and is itself
KLS-strength.  Thus (17.12) is a product theorem, not the missing
irreducible ANOVA inequality.

\subsection{Why the evident reverse arguments fail}

Write \(\mathcal E_\eta(u)=\int|\nabla u|^2\,d\eta\).  If \(f\) is an
attained normalized bottom mode of \(\mu\), define
\[
 H=\frac{f(X)+f(Y)}{\sqrt2},\qquad
 g=\E[H\mid S],\qquad r=H-g.
\]
Orthogonal projection and the weak eigenvalue equation give
\[
 \boxed{\qquad
 \mathcal E_\nu(g)
 =a\bigl(1-2\norm r_2^2\bigr)+\mathcal E_{\mu^2}(r).
 \qquad}                                              \tag{10.4}
\]
Conditional projection need not decrease energy.  For
\(\mu=\operatorname{Unif}[-1,1]\),
\[
 f(x)=\sqrt2\sin(\pi x/2),\qquad a=\pi^2/4.
\]
Writing \(t=\sqrt2s\), direct averaging on the sum fiber gives
\[
 g(s)=\frac4\pi\frac{\sin(\pi t/2)}{2-|t|},\qquad |t|<2.
\]
If \(J=\int_0^\pi\sin^2(v)\,dv/v\), then
\[
 \Var_\nu(g)=\frac{8J}{\pi^2},\qquad
 \mathcal E_\nu(g)=4(J-\tfrac12)>\frac{\pi^2}{4}
                  =\mathcal E_{\mu^2}(H),            \tag{10.5}
\]
where
\[
 J\ge\frac{\pi^2}{8}-\frac{\pi^4}{192}
 +\frac{\pi^6}{13824}+\frac5{16}+\frac1{4\pi^2}
 >\frac12+\frac{\pi^2}{16}.                          \tag{10.6}
\]
Thus discarding the fiber residual cannot reverse (SC).

Several scalar Lyapunov budgets fail for independent reasons.  For
\[
 X_k=\frac{\Gamma(k,1)-k}{\sqrt k},
\]
normalized self-convolution maps \(X_k\) to \(X_{2k}\), while
\[
                         \lambda_1(X_k)=\frac{k^2}{(k+1)^2}. \tag{10.7}
\]
Raw Fisher information is \(k/(k-2)\) for \(k>2\) and is infinite at the
mandatory hard-support cases \(k\le2\).  Gaussian regularizations of the
shifted exponential retain the limiting gap change
\(1/4\mapsto4/9\), while the decrement of every bounded inverse-Fisher
transform tends to zero.

There is also a smooth high-chaos obstruction.  Let \(H_m\) be the
probabilists' Hermite polynomial, fix even \(m\ge4\), and
variance-normalize the strongly log-concave density proportional to
\[
 \exp[-x^2/2-\varepsilon H_m(x)/\sqrt{m!}].
\]
For fixed \(m\) and \(\varepsilon\downarrow0\), its gap \(a\) and the gap
\(b\) after normalized self-convolution satisfy
\[
 a=1-\frac{m}{m-2}\varepsilon^2+O_m(\varepsilon^3),
 \quad
 b=1-\frac{m}{m-2}2^{2-m}\varepsilon^2
        +O_m(\varepsilon^3).                         \tag{10.8}
\]
For isotropic \(\eta=\mathcal L(Z)\), define
\[
 \mathcal F(\eta)
 =I\left(\frac{Z+G}{\sqrt2}\,\middle\|\,\gamma\right)
 =1-2\operatorname{mmse}(Z\mid Z+G).
\]
Then
\[
 \lim_{\varepsilon\downarrow0}
 \frac{\log(b/a)}
 {\mathcal F(\mu)-\mathcal F(\mathsf T\mu)}
 =\frac{2^m}{m-2}.                                   \tag{10.9}
\]
Letting \(m\) increase rules out a universal fixed-noise scalar charge.
Entropy survives this high-chaos test, but raw entropy and dependence
deficits are additive under products, normalization by dimension fails
under Gaussian padding, and bounded transforms saturate.  The correct
Hadamard sign is
\[
                         I(S;D)=h(S)+h(D)-2h(X);
\]
for shifted exponentials the value is Euler's constant.  These tests leave
open directional, operator-valued, and function-specific reverse
quantities.

\subsection{The Gaussian-channel affine gate}

Let \(Y=X+G\), let \(q\) be the law of \(Y\), and suppose that a smooth
\(g\) satisfies
\[
 \E_\mu g=0,\qquad \E_\mu[Xg]=0,\qquad
 h(y)=\E[g(X)\mid Y=y].
\]
The output score and affine orthogonality give
\[
                         |\E_q\nabla h|\le2^{-1/2}\norm g_2. \tag{10.10}
\]
For the posterior \(\mu_y\), write
\[
 A_y=\Cov_y(X),\qquad a_y=\E_y\nabla g,
\]
and let \(r_y\) be the residual after subtracting the posterior mean and
best affine predictor from \(g\).  Posterior strong convexity yields
\[
 \E_\mu\nabla g-\E_q\nabla h
 =\E[(I-A_Y)a_Y]-\E\Cov_Y(r_Y,X),                    \tag{10.11}
\]
\[
 \left|\E\Cov_Y(r_Y,X)\right|
 \le\norm{D^2g}_{L^2(\mu)}.                          \tag{10.12}
\]
The exact uncontrolled term is
\[
 \boxed{\qquad
 \left|\E[(I-A_Y)a_Y]\right|
 \le C\norm{D^2g}_{L^2(\mu)}.
 \qquad}                                              \tag{AFF}
\]
A pointwise posterior form is false, since a posterior-affine function has
zero Hessian.  The integrated estimate (AFF) remains unproved and is a
KLS-strength near-linear gate; inserting a dimension-free Poincar\'e
inequality for \(q\) would be circular.

\subsection{The exact weighted-Fisher ledger}

In the regular conditional-slice setting, let
\[
 r(s,z)=e^{-U(s,z)},\quad
 \rho(s)=\int r(s,z)\,dz=e^{-\Phi(s)},\quad
 q_s(z)=\frac{r(s,z)}{\rho(s)}.
\]
With
\[
 L_s=\Delta_z-\nabla_zU(s,\cdot)\cdot\nabla_z,\quad
 L_sg_s=\ell_s:=\partial_s\log q_s,\quad
 F_s=\nabla_zg_s,\quad H_s=D^2_{zz}U,
\]
put
\[
 C_s=\E_s\norm{D_zF_s}_{\HS}^2,\qquad
 B_s=\E_s\ip{H_sF_s}{F_s}.
\]
Conditional Bochner gives
\[
                         I_\perp(s):=\E_s\ell_s^2=C_s+B_s. \tag{10.13}
\]
Define
\[
 \alpha_s=\partial_sg_s-\tfrac12|F_s|^2,\qquad
 \mathcal Q_s=D^2U[(1,-F_s),(1,-F_s)].
\]
Joint convexity makes \(\mathcal Q_s\ge0\), and direct differentiation
gives
\[
 L_s\alpha_s=\Phi''-\mathcal Q_s-\norm{D_zF_s}_{\HS}^2,
 \qquad
 \Phi''=\E_s\mathcal Q_s+C_s.                        \tag{10.14}
\]
If the scalar marginal is centered and variance one and \(\tau\) is its
canonical Stein kernel, then
\[
 \boxed{\quad
 \int\rho\tau^2(C_s+\E_s\mathcal Q_s)\,ds
 =\int\rho\tau^2\Phi''\,ds
 =1-\E_\rho(\tau')^2\le1.
 \quad}                                               \tag{10.15}
\]
Thus deformation and material curvature already have a universal budget.
The conditional curvature term \(B_s\) does not: in the full-space Hodge
identity its apparent contribution cancels exactly against the horizontal
and mixed terms.

For the conditional centroid \(m(s)=\E_sZ\), the same calculation gives
\[
 m'=-\E_sF_s,\quad
 m''=-\E_s[(Z-m)(\mathcal Q_s+\norm{D_zF_s}_{\HS}^2)],
 \quad \E[\tau(S)m'(S)]=0.                           \tag{10.16}
\]
These global-in-\(s\) identities are indispensable: a slice-local
replacement is false below.

\subsection{Two complete weighted-Fisher subclasses}

\paragraph{Pure translations.}
Suppose before transverse smoothing that
\[
 p_s(x)=Z_\varphi^{-1}e^{-\varphi(x-m(s))}
\]
and that the joint law is isotropic and log-concave.  If \(Y\) has density
proportional to \(e^{-\varphi}\), isotropy gives \(\Cov(Y)\preceq I\).
Including subgradients and normal cones in the nonsmooth case,
\[
 \overline{\operatorname{conv}}\{\partial\varphi(y)\}
 \supset\frac3{16}B_2^d.                             \tag{10.17}
\]
Testing joint convexity against this score-range inradius yields,
in distributions,
\[
                         |Dm'|\le\frac{16}{3}D\Phi'. \tag{10.18}
\]
Cross-isotropy and one-dimensional Stein estimates give
\[
 \E[\tau^2|m'|^2]
 \le3208\left(\frac{16}{3}\right)^2.                 \tag{10.19}
\]
After unit transverse Gaussian smoothing, Gaussian-channel contraction
gives \(I_\perp(s)\le|m'(s)|^2\).  Consequently
\[
                         \int\rho\tau^2I_\perp\,ds\le C \tag{10.20}
\]
for every pure translating family, uniformly in transverse dimension.
The finite-difference proof includes asymmetric fibers and hard endpoints.

\paragraph{Smooth post-noise Gaussian fibers.}
Consider
\[
 p(s,z)=Z^{-1}\exp\left[
 -W(s)-\tfrac12(z-m(s))^TQ(s)(z-m(s))\right],
 \qquad R(s)=Q(s)^{-1}.
\]
Assume \(W,m,R\in C^2\), \(R\succ0\), the joint potential is convex, the
\(S\)-marginal is centered with variance one, and
\[
 \E[S\,m(S)]=0,\qquad
 I\preceq R(s)\ \hbox{for every \(s\)},\qquad
 \E R(S)\preceq\Lambda I.                            \tag{10.21}
\]
These are precisely the smooth full-support hypotheses used here; the
lower bound on \(R\) is the unit transverse noise.

Put \(y=z-m(s)\).  The exact centered Poisson field is
\[
 F_s=-m'(s)+A_sy,\qquad A_sR+RA_s=-R'.
\]
Writing
\[
 B_s=m'(s)^TR^{-1}m'(s)+B_s^0,\qquad
 J_s=\frac12\Tr(R^{-1}R'R^{-1}R'),
\]
diagonalization of \(R(s)\) at fixed \(s\), followed by pairing the
\((i,j)\) and \((j,i)\) terms, gives the noncommutative identity
\[
                         \boxed{C_s+B_s^0=J_s}.       \tag{10.22}
\]
No derivative of the diagonalizing frame is taken.  Joint convexity gives
\[
 \Phi''
 =W''+\frac12\Tr(R^{-1}(-R''))+J_s\ge J_s,           \tag{10.23}
\]
and hence
\[
                         \int\rho\tau^2(C_s+B_s^0)\,ds\le1. \tag{10.24}
\]

The centroid theorem has the following audited scope: for a \(C^2\),
full-support jointly convex Gaussian family satisfying the marginal and
mixed-isotropy hypotheses above, but requiring only
\(\E R\preceq\Lambda I\), matrix concavity and the Schur condition imply
\[
                         \E[\tau^2|m'|^2]\le C_{\rm cent}\Lambda. \tag{10.25}
\]
This allows noncommuting, rotating covariance matrices.  Under the
post-noise lower bound \(R\succeq I\), one has \(R^{-1}\preceq I\), so
\[
 \boxed{\qquad
 \int\rho(s)\tau(s)^2(C_s+B_s)\,ds
 \le1+C_{\rm cent}\Lambda.
 \qquad}                                              \tag{10.26}
\]
No nonsmooth Gaussian limit and no zero-noise limit is asserted.  A
separate hard-interval calculation controls geometric centroid velocity,
including wedge curvature atoms, but its identification with the
post-noise Poisson field has not been proved.

\subsection{A cone forcing global Stein weighting}

For \(a\ge2\), let
\[
 q_a=\operatorname{Unif}[-a,a]*\gamma,\qquad U_a=-\log q_a.
\]
The mean-zero gradient field \(F_a(x)=x/a\) satisfies
\[
 \E_{q_a}|F_a'|^2=a^{-2},\qquad
 \E_{q_a}[U_a''F_a^2]\ge\frac1{64a}.                 \tag{10.27}
\]
More decisively, the same divergence occurs for the exact Poisson field
of one fixed isotropic law.  Start from
\[
 p(r,x)=\tfrac12e^{-r}\1_{\{r>0,\ |x|<r\}},\qquad
 S=\frac{R-2}{\sqrt2},\qquad Y=\frac X{\sqrt2}.
\]
Then \(Y\mid S=s\) is uniform on \([-a,a]\), where
\(a=s+\sqrt2>0\).  After unit Gaussian smoothing in \(Y\), its exact
centered Poisson field obeys
\[
 \E_sF_s=0,\qquad C_s\le a^{-2},\qquad
 B_s\ge\frac1{256a}.                                 \tag{10.28}
\]
Thus no universal slice-by-slice estimate
\(B_s\lesssim C_s+|\E_sF_s|^2\) can hold.  Nevertheless
\[
 \rho(s)=2ae^{-\sqrt2a}\1_{\{a>0\}},\qquad
 \tau(s)=\frac a{\sqrt2},\qquad B_s\le1,
\]
and therefore
\[
                         \int\rho\tau^2B_s\,ds\le\frac32. \tag{10.29}
\]
The obstruction is entirely in the centered shape residual, not centroid
motion.  It forces a proof to retain the joint \(s\)-geometry and Stein
weight.

\subsection{Exact remaining gates}

The weighted route would close its general slice estimate if one proved
\[
 \boxed{\qquad
 \int\rho\tau^2B_s\,ds
 \le C\int\rho\tau^2(C_s+\E_s\mathcal Q_s)\,ds.
 \qquad}                                              \tag{WK}
\]
The cone proves that (WK) must be integrated in \(s\), while the
translation theorem removes pure centroid motion from the unresolved
sector.  What remains is genuinely nonlinear shape change.  A proposed
scalar shape-acceleration shortcut is invalid: midpoint expansion cancels
the \(U_zm''\) term and yields only affine-path curvature.

The two surviving near-linear gates are (AFF) and (WK), but even a full
weighted-Fisher or directional-MMSE theorem would currently close only the
near-linear spectral branch.  The tensor-extremizer program still lacks a
proved dimension-free dichotomy forcing every hypothetical small-gap mode
to be near-linear, or an alternative argument applying the weighted
estimate to every low mode.  Neither (AFF), (WK), nor this global
nonlinear/near-linear dichotomy is assumed here.

\section{Mechanisms ruled out by explicit tests}

The negative results are as useful as the positive ones because they
prevent equivalent forms of the conjecture from being reintroduced under
new notation.

\begin{longtable}{>{\raggedright\arraybackslash}p{0.25\textwidth}
                  >{\raggedright\arraybackslash}p{0.68\textwidth}}
\toprule
Proposed mechanism & Audited obstruction \\
\midrule
\endfirsthead
\toprule
Proposed mechanism & Audited obstruction \\
\midrule
\endhead
Thin shell alone controls (TM) & Projection traces permit the harmonic
spectrum \(j^{-1/2}\), leaving \(\sqrt{\log n}\). \\
\addlinespace
Full-gradient interpolation & The first Neumann mode of the isotropic ball
violates the natural multiplicative interpolation by \(\sqrt n\). \\
\addlinespace
Reverse conditional projection & For the first interval mode,
\(\mathcal E_\nu(\E[H\mid S])>\mathcal E_{\mu^2}(H)\); the exact residual
identity (10.4) has no favorable sign. \\
\addlinespace
Scalar reverse-convolution Lyapunov & Hard-support Gamma laws defeat
bounded inverse-Fisher budgets, high Hermite modes defeat fixed-noise
Fisher/MMSE charging, and additive entropy budgets fail under products or
Gaussian padding. \\
\addlinespace
Trace localization \(\Rightarrow\) selected direction & A one-spike
product-integral path has \(O(n)\) total trace but polylogarithmic
amplification in the chosen direction. \\
\addlinespace
Coarea or determinant control of (FA) & The flow expectation is at a fixed
initial tilt.  A linear map \(\operatorname{diag}(\sqrt n,1,\ldots,1)\)
has acceptable trace and determinant scale but normal stretch \(n\). \\
\addlinespace
Strong self-concordance \(\Rightarrow\) KLS & On an isotropic cube, an
explicit strongly self-concordant barrier has weighted Poincar\'e constant
of order \(n\).  Barrier axioms give no covariance comparison. \\
\addlinespace
Bakry--\'Emery for the canonical metric & For the direct Cheng--Yau
potential the natural curvature has the wrong sign; for uniform measure
the curvature-to-metric ratio already tends to zero at an interval
boundary. \\
\addlinespace
Pointwise barrier-metric comparison & It fails on the simplex, skew cones,
and product exponentials; large inverse metric can occur on rare sets. \\
\addlinespace
Tensor multiplicity & Exact splitting never exceeds the convex buffer, and
the low modes remain inside the allowed small-rank localization tail. \\
\addlinespace
Interior Brascamp--Lieb for slices & Uniform intervals and centered
exponentials carry covariance motion entirely through moving-boundary
terms. \\
\addlinespace
Pointwise post-noise curvature--Korn & The exact cone Poisson field has
\(C_s\le a^{-2}\), \(B_s\ge(256a)^{-1}\), and zero mean, although its
Stein-weighted integral is at most \(3/2\). \\
\addlinespace
Scalar shape acceleration & Midpoint expansion cancels the proposed
\(U_zm''\) term and leaves only affine-path curvature. \\
\addlinespace
Finite Gaussian chaos & Exponential examples have factorial mixed moments;
the relevant Hermite expansion can have zero radius of convergence. \\
\addlinespace
Entropy-controlled convex certificate & A one-dimensional zero-deficit
transport defeats every prescribed-orientation signed convex certificate. \\
\bottomrule
\end{longtable}

\section{Prioritized future directions}

The following program is ordered by how sharply each step is formulated
and how independently it can be tested.

\subsection{Priority I: posterior covariance and low-mode alignment}

\begin{enumerate}
\item Prove or disprove (DMMSE).  First establish the proposed
dimension-descent estimate (7.9) with all conditional derivatives and tail
bounds written rigorously.  Then identify an additional rigidity input
that removes its logarithm.  Candidate inputs include a tail theorem for
the rank-one survivor variance or a nonlinear-information bound forced by
joint log-concavity.

\item Attack (FA) directly as a matrix-martingale statement.  Any proof
must use that \(c_t\) is a posterior gradient of a fixed scalar low mode;
trace bounds and arbitrary terminal directions are insufficient.  A useful
intermediate target is a bound on amplification-weighted occupation of
\((A_t-3I)_+\) along the direction selected by \(c_t\).

\item Relate the two targets.  The posterior split (7.2) shows that
(DMMSE) closes very-near-linear modes, while (FA) is a pathwise refinement.
A proof should quantify precisely how much residual size \(\delta\) each
one tolerates.

\item Prove or disprove (AFF).  Equations (10.11)--(10.12) already control
every posterior-nonlinear component; any proof must control the remaining
posterior-affine term globally rather than by a false pointwise Hessian
estimate.
\end{enumerate}

\subsection{Priority II: non-Gaussian slice-curvature compensation}

The smooth post-noise Gaussian theorem proves the full weighted-Fisher
bound by the exact identity \(C_s+B_s^0=J_s\), even for noncommuting
covariance motion.  The next target is the integrated estimate (WK) for
general slices.  The cone calculation forbids replacing it by a
slice-by-slice Korn inequality, while the translation theorem shows that
pure centroid motion is no longer the issue.

One possible route to (WK) is a boundary-inclusive version of the
Brascamp--Lieb stability estimate
\[
 \norm{C^{-1/2}\Cov_t(g,(Y-m)(Y-m)^T)C^{-1/2}}_{\HS}^2
 \lesssim \mathcal D_{\mathrm{BL}}(g)
 +\text{boundary/Reilly charge}.                     \tag{9.1}
\]
The endpoint term must reproduce the interval and exponential examples
exactly.  Before attempting full generality, test (9.1) on polytopal slices
with rotating facet normals and on smooth approximations of boxes.

The affine-Hessian inequality (AH) is a clean alternative formulation.
It should first be proved or falsified for centrally symmetric but
dependent bodies, uniform polytopes, and skew cones.  Any proof for all
gradients would contain a KLS-strength near-linear statement, so a
score-specific version may be more realistic.

\subsection{Priority III: exact canonical metric, only for low modes}

Strong self-concordance alone is blocked.  A narrower question remains:
for the exact direct Cheng--Yau metric \(H=D^2F\), can a first Euclidean
eigenfunction satisfy
\[
 \int\ip{H^{-1}\nabla f}{\nabla f}\dd\mu
 \le C\left(
 \int|\nabla f|^2\dd\mu+
 \int\norm{D^2f}_{\HS}^2\dd\mu\right)?              \tag{9.2}
\]
Bochner supplies the second term at order \(\lambda^2\).  The central
issue is whether a low mode can concentrate its gradient on the rare sets
where \(H^{-1}\) is large.  The simplex and skew-cone calculations provide
sharp stress tests.  This direction is worthwhile only if it uses the
exact canonical equation; generic barrier axioms have already been ruled
out.

\subsection{Priority IV: extremal structure}

If one pursues an extremal sequence, impose all known constraints at once:

\begin{itemize}
\item witnesses are nonlinear and non-radial;
\item tensor-product explanations are factored out;
\item near-linear witnesses must satisfy the posterior and localization
constraints (6.5), (FA), and (7.5);
\item any tensor stationarity argument must pay the exact convex-buffer
cost (8.3);
\item any slicing argument must retain moving-boundary curvature.
\end{itemize}

A useful structure theorem would state that every counterexample sequence
has one of two mutually exclusive features: a direction with vanishing
unit-noise posterior MMSE, or a genuinely high-rank noncommuting
conditional-covariance motion.  The two active programs above are designed
to exclude precisely those alternatives.

\section{Audit requirements for any future proof}

Before a candidate argument is regarded as a proof, it should pass the
following checks.

\begin{enumerate}
\item Track \(n\), smoothing, truncation, localization time, stopping rank,
and bootstrap depth symbolically.  The final constant must be independent
of all of them.
\item Verify stochastic existence, stopping, integrability, and true
martingale properties rather than formal local-martingale calculations.
\item Pass from smooth, compact, strongly log-concave approximations to
arbitrary log-concave measures without changing the constant.
\item Treat lower-dimensional affine supports intrinsically and exclude
only point masses.
\item Cover every locally Lipschitz \(L^2\) test function in a Poincar\'e
proof.
\item Scan for circular assumptions: uniform Poincar\'e, (M), (TM), (MG),
dimension-free survivor covariance, (AFF), (WK), and strong
self-concordance of the entropic barrier are all KLS-strength gates unless
independently proved.
\item Test the constant chain on the cube, regular simplex, crosspolytope,
product exponential, Euclidean ball, and a skew non-symmetric cone.
\item Have a clean-room reader reprove every load-bearing lemma from its
statement alone.
\end{enumerate}

\section{Conclusion}

The investigation has not resolved KLS.  The strict convolution inequality
(SC), including the unequal-input form (USC), remains a forward
amplification statement.  The independent posterior comparison
\[
 a\ge\frac{b}{2+b}
\]
does give a dimension-free reduction to positive analytic isotropic
outputs with \(0\preceq D^2\widetilde U\preceq2I\).  Along a hypothetical
small-gap sequence this comparison is now known to saturate to second
order: the unscaled output gap is \(a+O(a^2)\), and bottom input modes
saturate the posterior Poincar\'e and covariance inequalities on average.

Checkpoint 17 strengthens the local rigidity picture without closing its
global gate.  The Hilbert--Schmidt posterior theorem removes the former
trace loss, and strongly-log-concave affine stability shows that bottom
modes are posterior-affine on average up to squared error \(O(a^2)\).
Under a uniform simple posterior eigengap, the low spectral projection has
dimension-free derivative energy.  Nevertheless the twice-Bochner ledger
still contains an \(O_\kappa(\lambda)\) transverse saturation remainder
and a signed longitudinal quotient term.  The longitudinal term is
coercive with constant \(1/24\) after the line is fixed, but explicit
Gaussian examples rule out generic amplitude-weighted connectivity as a
way to fix it.

The exact Gaussian-jet identity also sharpens the nonlinear branch.  It
proves the affine-interaction estimate with constant \(900\) for arbitrary
products of one-dimensional isotropic log-concave laws.  Extending that
argument to irreducible laws would require a new operator-level
matrix-Stein or equivalent coercivity mechanism; known trace budgets and
the input Poincar\'e inequality are insufficient.

Accordingly, the current frontier has two coupled components: global
posterior direction and amplitude synchronization strong enough to absorb
simple-eigenspace remainders, and a dimension-free affine-interaction or
nonlinear/near-linear dichotomy for general log-concave laws.  The older
weighted-Fisher gates (AFF), (WK), (DMMSE), and (FA) remain alternative
formulations, not assumptions.  None of the audited reductions, stability
theorems, fixed-line arguments, or subclass closures in this note is a
complete proof of KLS.

\begin{thebibliography}{99}

\bibitem{Borell1975}
C.~Borell,
\newblock Convex set functions in \(d\)-space,
\newblock \emph{Periodica Mathematica Hungarica} 6 (1975), 111--136.

\bibitem{KLS1995}
R.~Kannan, L.~Lov\'asz, and M.~Simonovits,
\newblock Isoperimetric problems for convex bodies and a localization lemma,
\newblock \emph{Discrete \& Computational Geometry} 13 (1995), 541--559.

\bibitem{Ledoux2001}
M.~Ledoux,
\newblock \emph{The Concentration of Measure Phenomenon},
\newblock American Mathematical Society, 2001.

\bibitem{Milman2009}
E.~Milman,
\newblock On the role of convexity in isoperimetry, spectral gap and
concentration,
\newblock \emph{Inventiones Mathematicae} 177 (2009), 1--43.

\bibitem{Eldan2013}
R.~Eldan,
\newblock Thin shell implies spectral gap up to polylog via a stochastic
localization scheme,
\newblock \emph{Geometric and Functional Analysis} 23 (2013), 532--569.

\bibitem{BubeckEldan2015}
S.~Bubeck and R.~Eldan,
\newblock The entropic barrier: a simple and optimal universal
self-concordant barrier,
\newblock in \emph{Proceedings of COLT 2015}, 279--279.

\bibitem{Hildebrand2014}
R.~Hildebrand,
\newblock Canonical barriers on convex cones,
\newblock \emph{Mathematics of Operations Research} 39 (2014), 841--850.

\bibitem{LaddhaLeeVempala2020}
A.~Laddha, Y.~T.~Lee, and S.~Vempala,
\newblock Strong self-concordance and sampling,
\newblock in \emph{Proceedings of STOC 2020}, 1212--1222.

\end{thebibliography}

\end{document}
